\subsection{Physics Constraints Analysis}
\label{sec:discussion-physics}

Physics constraints embed domain knowledge by encoding the governing physics of Lamb wave propagation directly into the training objective. Unlike a generic MSE loss that treats the signal as a pixel-level regression problem, the three physics terms in $\mathcal{L}_{\text{physics}}$ (Section~\ref{sec:physics-loss}) penalize deviations in diagnostically specific acoustic features: arrival time, peak amplitude, and waveform morphology. Arrival time preservation penalizes errors in first-break detection, directly impacting the accuracy of defect localization through time-of-flight measurement. Amplitude preservation penalizes distortion in peak signal amplitudes, which is critical for defect sizing since acoustic energy reflectivity at a defect interface scales with flaw severity. Waveform morphology is enforced via Lin's concordance correlation coefficient applied to the signal envelope, penalizing shape distortion that would otherwise obscure modal conversion effects associated with complex defect geometries.

MSE fails for this application because it optimizes aggregate waveform similarity, implicitly favoring overly smooth outputs that suppress high-frequency transients including arrival-time discontinuities. The ablation results in Section~\ref{sec:results-ablation} confirm this mechanism: removing $\mathcal{L}_{\text{physics}}$ while retaining competitive SNR ($12.1$ vs $12.4$ dB) causes TOA error to increase from $0.42$ to $0.68$ $\mu$s, amplitude error to rise from $3.1\%$ to $7.4\%$, and F1 score to drop from $0.91$ to $0.79$. The physics constraints thus function as an inductive bias that guides the network toward physically plausible solutions without degrading denoising capacity.

A key practical consequence is generalization benefit. Because the physics terms encode universal acoustic wave physics rather than specimen-specific signal characteristics, the learned feature preservation behavior transfers across specimen geometries. This explains the intact-to-defect and cross-thickness generalization results reported in Section~\ref{sec:results-generalization}, where PMDF-Net with physics loss retained near-intact performance on defect and thickness-shifted test sets.
